\begin{table}[htbp]
\centering
\caption{Order-book anatomy of price clustering}
\label{tab:rq3}
\small
\begin{threeparttable}
\begin{tabular}{@{}lrrrrrrr@{}}
\toprule
Dependent variable & $L^{0}$ & SE & $L^{0C}$ & SE & RDepth & SE & Stock-days \\
\midrule
$M^{BLarge0}$ & 0.0029 & (0.0524) & -0.0092 & (0.0136) & 0.2565*** & (0.0222) & 9,527 \\
$\Delta\mathit{Imp}^{BLarge}$ & 3.6734 & (9.3067) & -0.7720 & (2.0219) & 17.6809* & (10.0115) & 5,366 \\
$M^{SLarge0}$ & 0.1501 & (0.0983) & 0.0117 & (0.0197) & 0.1763*** & (0.0408) & 9,537 \\
$\Delta\mathit{Imp}^{SLarge}$ & 6.0726 & (7.7144) & -2.4585 & (2.1336) & 4.1140 & (4.0574) & 5,722 \\
\bottomrule
\end{tabular}
\begin{tablenotes}[flushleft]\footnotesize
\item Analogues of Ohta's (2026) equations (1) and (2), with book-derived variables replacing the margin-trading and ownership proxies, which are not available here. $L^{0}$ is the round-price share of limit-order volume submitted within the visible book, $L^{0C}$ the ratio of cancelled to submitted volume at round prices, and RDepth the round-price share of visible ten-level resting depth. Sell-side book variables are paired with buy-initiated outcomes, since a buy market order executes against sell limit orders. All regressors are lagged one day. Stock and day fixed effects, standard errors clustered by stock and by day. The hypothesis predicts a positive coefficient on $L^{0}$ and RDepth and a negative one on $L^{0C}$. $^{*}p<0.1$, $^{**}p<0.05$, $^{***}p<0.01$.
\end{tablenotes}
\end{threeparttable}
\end{table}
